\chapter{贯穿项目 从两连杆到七轴相机云台}
\label{ch:integration}

这一章把前面的工具放到同一条链上。先在模型明确的两连杆对象上核对运动学、动力学、辨识和跟踪，再讨论七轴云台的底座、相机、柔顺与冗余。原稿记录的 NERO 回正问题作为诊断情景使用；没有真机日志支持的结论，不能在教材里写成已定位的原因。

\section{两连杆项目的统一参数和边界}

两连杆长度均为 $0.5\,\mathrm m$，名义质量均为 $1\,\mathrm{kg}$，质心位于半长处，绕运动平面法向的质心惯量为 $1/48\,\mathrm{kg\,m^2}$；重力沿 $-y$，第一关节角 $q_1$ 从水平正方向起算，$q_2$ 是相对第一杆的角度。起点 $q_A=(0,\pi/2)$，终点 $q_B=(\pi/3,\pi/3)$，三秒五次轨迹，之后保持一秒。

本章实际生成的仿真采用 $200\,\mathrm{Hz}$ 控制、$0.5\,\mathrm{ms}$ 数值积分、一个控制周期的命令延迟和 $4\,\mathrm{ms}$ 执行器时间常数。对象质量设为 $(1.1,0.9)\,\mathrm{kg}$，杆惯量按质量同比缩放；摩擦为平滑库仑加黏性模型，参数与辨识数据一起记录。所有这些是\textbf{可复现的教学设定，不是 NERO 规格或测量}。

\subsection{先得到不依赖任何库的核对点}

对名义模型，在起点，末端 $p_A=(0.5,0.5)\,\mathrm m$，
\begin{equation}
J_A=\begin{bmatrix}-0.5&-0.5\\0.5&0\end{bmatrix},\quad
M_A=\begin{bmatrix}5/12&1/12\\1/12&1/12\end{bmatrix},\quad
g_A=(7.3575,0)^\mathsf T.
\end{equation}
若要瞬时末端速度 $(0,0.1)\,\mathrm{m/s}$，第二行给 $0.5\dot q_1=0.1$，所以 $\dot q_1=0.2$；第一行要求 $\dot q_1+\dot q_2=0$，故 $\dot q_2=-0.2\,\mathrm{rad/s}$。

若施加环境力 $f=(2,-3)\,\mathrm N$，其关节外力矩为
\[
J_A^\mathsf Tf=(-2.5,-1)^\mathsf T\,\mathrm{N\,m}.
\]
虚功核对：$\dot q^\mathsf TJ^\mathsf Tf=0.2(-2.5)+(-0.2)(-1)=-0.3\,\mathrm W$，同时 $v^\mathsf Tf=(0,0.1)\cdot(2,-3)=-0.3\,\mathrm W$。两边相等，可以同时检查关节顺序、转置和力方向。

\section{如何让辨识结果进入控制器}

两连杆刚体动力学可写成五个合成参数：
\begin{align}
\beta_1&=I_1+I_2+m_1c_1^2+m_2(l_1^2+c_2^2),\\
\beta_2&=m_2l_1c_2,\qquad \beta_3=I_2+m_2c_2^2,\\
\beta_4&=m_1c_1+m_2l_1,\qquad \beta_5=m_2c_2.
\end{align}
前三个有惯量单位，后两个为质量一阶矩。它们不是七轴机器人“43个基参数”的缩小版计数规则，而是当前 2R 方程的一种线性参数化。由于本例质心沿杆轴、$l_1$ 已知，存在 $\beta_2=l_1\beta_5$；若再采用已知的匀质杆关系 $I_i=m_il_i^2/12$ 和 $c_i=l_i/2$，五个组合全部由两个质量决定。因此本章实际实验只辨识两个质量，不将五个组合当成五个独立物理未知量。一般机器人完整辨识仍需采用第8章的秩分析与物理约束。

例如将第一行回归特征写为 $Y_1=(y_{11},y_{12},y_{13},y_{14},y_{15})$，逐项为
\begin{align}
y_{11}&=\ddot q_1,\qquad y_{13}=\ddot q_2,\\
y_{12}&=2\cos q_2\ddot q_1+\cos q_2\ddot q_2
-\sin q_2(2\dot q_1\dot q_2+\dot q_2^2),\\
y_{14}&=g_0\cos q_1,\qquad y_{15}=g_0\cos(q_1+q_2).
\end{align}
这里的余弦和正弦均显式写出，避免与质心距离 $c_2$ 混用。

教材配套计算先用一条解析可微的合成激励轨迹产生位置、速度和加速度，再由对象真值生成带固定随机种子噪声的力矩记录，利用已知匀质杆几何拟合两个质量和四个摩擦系数，再重建上述合成参数；另一条轨迹用于验证。控制器只使用拟合结果或名义模型，仿真对象单独保留真值。这里为了隔离辨识结构问题，位置与导数取解析值；真实编码器的噪声处理在第9章与第21章讨论，不能把本实验误认为完整复现了真机采集链。

\section{四种控制器的实际数值比较}

使用相同的关节反馈增益 $K_p=\operatorname{diag}(30,12)$、$K_d=\operatorname{diag}(8,3)$，分别计算：纯 PD；名义重力补偿加 PD；名义逆动力学轨迹前馈加 PD；辨识模型及摩擦前馈加 PD。这里的反馈增益分别是力矩每弧度和力矩每角速度的量纲。

逆动力学前馈使用参考状态：
\begin{equation}
\tau_{\rm ff}=\hat M(q_d)\ddot q_d+\hat C(q_d,\dot q_d)\dot q_d+\hat g(q_d)+\hat\tau_f(\dot q_d),
\end{equation}
再加 $-K_p(q-q_d)-K_d(\dot q-\dot q_d)$。这与把反馈加速度也乘以实测状态惯量的计算力矩控制有区别。前馈中没有使用实测加速度，也不需要通过二阶差分去制造一条高噪声信号。

\begin{figure}[htbp]\centering
\includegraphics[width=0.97\linewidth]{figures/rt_integrated_2r.pdf}
\caption{四种控制器在同一教学对象、同一轨迹、相同延迟和反馈增益下的数值结果。末端误差按真实对象的几何计算；不存在几何参数误差，因此本图只比较动态跟踪链。}
\end{figure}

\begin{table}[htbp]\centering\small
\begin{tabularx}{\linewidth}{Xrrrr}\toprule
控制器 & RMS / mm & 最大 / mm & 4 s末值 / mm & 饱和拍数\\\midrule
纯 PD & 140.190 & 212.088 & 56.603 & 0 \\
名义重力补偿加 PD & 3.052 & 5.745 & 3.861 & 0 \\
名义动力学前馈加 PD & 5.729 & 7.640 & 6.890 & 0 \\
辨识模型与摩擦前馈加 PD & 1.414 & 3.952 & 0.117 & 0 \\
\bottomrule\end{tabularx}
\caption{统一两连杆实验的实际计算结果。误差在200 Hz的0至4 s共801个记录点上统计，末值为4 s、到达后保持1 s的误差。饱和拍数统计4 s之前800次控制计算中至少一轴指令被裁剪的次数。}\end{table}

RMS为末端二维误差向量模长的均方根；最大值也是上述记录点上的最大误差，不是逐坐标百分比误差。

训练轨迹包含 3000 个时刻，每时刻两条力矩方程。独立验证的两轴力矩 RMSE 分别为 $0.02032$ 和 $0.01992\,\mathrm{N\,m}$；验证力矩含标准差 $0.02\,\mathrm{N\,m}$ 的独立合成噪声。

本实验将匀质杆几何和摩擦平滑速度尺度 $v_s=0.02\,\mathrm{rad/s}$ 作为已知设定，只辨识两个质量与四个摩擦系数，共六个参数。估计值为
\[
\begin{aligned}
(\hat m_1,\hat m_2)&=(1.09992,0.89993)\,\mathrm{kg},\\
(\hat F_{c1},\hat F_{c2})&=(0.14999,0.10173)\,\mathrm{N\,m},\\
(\hat b_1,\hat b_2)&=(0.07985,0.04637)\,\mathrm{N\,m\,s/rad}.
\end{aligned}
\]

关节力矩上限为 $(15,6)\,\mathrm{N\,m}$；四组实验的延迟队列与执行器力矩从零初始化，所以起始短暂下沉也计入误差。该初始化用于比较延迟影响，不代表推荐的真实硬件激活流程。



表中数据由配套脚本实际计算。注意，名义动力学前馈的 RMS 为 $5.729\,\mathrm{mm}$，反而大于名义重力补偿的 $3.052\,\mathrm{mm}$。两者并不是“在同一个控制律上准确地多加一项”：重力补偿在实测 $q$ 计算，而轨迹前馈在参考 $q_d$ 计算，模型质量有偏差，摩擦又未补偿，时延也使前馈与对象状态错位。增加模型项不自动保证误差下降。

不能拿某一列单独证明某算法普遍更好。改变轨迹、摩擦模型、滤波、饱和或时延，排序可能改变；本实验最有价值的是把“补了什么”和“误差改变多少”对应起来。所有曲线都有相同评价区间和单位，才能公平比较。

\subsection{补偿的求值点为什么改变局部刚度}

在保持目标 $q_d$ 附近，令 $G=\partial g/\partial q|_{q_d}$，$\hat G=\partial\hat g/\partial q|_{q_d}$。忽略高阶项，将重力展开为 $g(q_d+e)\approx g(q_d)+Ge$。若补偿固定为 $\hat g(q_d)$，恢复项近似为 $(K_p+G)e$；若补偿随实测位置变化为 $\hat g(q)$，则近似为 $(K_p+G-\hat G)e$。两种控制器对相同误差的局部恢复作用已经不同。

在本实验终点，这两个矩阵约为
\[
K_p+G=\begin{bmatrix}21.9291&-1.9115\\-1.9115&10.0885\end{bmatrix},\qquad
K_p+G-\hat G=\begin{bmatrix}30.4248&0.2124\\0.2124&12.2124\end{bmatrix}.
\]
单位为 $\mathrm{N\,m/rad}$。这只是保持阶段的一阶解释；完整轨迹还受惯量误差、摩擦残差和时延影响。若想单独评估“多加动力学项”的价值，应固定重力求值方式，逐项增加补偿并重新比较，避免把同时改变的机制归为一种原因。

\section{七轴云台首先是一个带柔性的支撑系统}

相机在世界中的位姿为
\begin{equation}
T_{WC}=T_{WB}T_{BE}(q)T_{EC}.
\end{equation}
$T_{WB}$ 是基座在世界中的位姿，将基座分量转换到世界系；$T_{BE}(q)$ 是工具在基座中的位姿，来自机械臂正运动学；$T_{EC}$ 是相机在工具中的固定安装位姿。关节编码器主要观测中间一段；底座移动、支架变形或相机安装松动，不会自动出现在关节位置误差里。

小角度基座扰动对相机位置的一阶影响包含
\begin{equation}
\delta p_{WC}\approx\delta p_{WB}
-R_{WB}[p_{BC}]_\times\delta\theta_B
+R_{WB}J_{p,C}\delta q.
\end{equation}
本式中 $\delta p_{WB}$ 在世界系表达，$\delta\theta_B$ 为在基座系表达的右扰动小角，$J_{p,C}$ 是相机参考点在基座系中的位置雅可比。因此机械臂关节保持得很好，也可能拍出抖动画面。若相机距基座 $0.8\,\mathrm m$，基座角度晃动 $0.1^\circ=0.0017453\,\mathrm{rad}$，垂直于力臂方向的位移约为 $0.8\times0.0017453=1.40\,\mathrm{mm}$。若相机等效焦距为 $1000$ 像素，对图像中心附近的特征，纯小旋转可造成约 $1.75$ 像素的角向画面位移；有限深度下的相机平移还会产生另一项图像运动。

\begin{figure}[htbp]\centering
\includegraphics[width=0.94\linewidth]{figures/rt_gimbal_base.pdf}
\caption{基座角度误差通过力臂放大成相机位移。右图为小角度一阶计算，不代表某种支架的实测刚度。控制器没有测到的基座运动，不能仅靠关节误差消除。}
\end{figure}

\subsection{从允许角度误差反推支撑刚度}

用一维转动弹簧近似底座：$\tau_b=k_\theta\theta_b$。若任务下变化力矩为 $20\,\mathrm{N\,m}$，允许静态角度变化为 $0.05^\circ$，则
\[
k_\theta\ge\frac{20}{0.05\pi/180}=22918\,\mathrm{N\,m/rad}.
\]
这是线性静态刚度估算。动态还需要等效转动惯量 $I_b$ 和阻尼 $d_b$，模型为 $I_b\ddot\theta_b+d_b\dot\theta_b+k_\theta\theta_b=\tau_b$，自然频率 $\omega_n=\sqrt{k_\theta/I_b}$。提高控制增益不一定能消除底座共振，可能反而把更多力矩变化传给支架。

\section{没到点至少有三种不同误差}

定义目标末端位置 $p_d$，规划出的关节目标 $q_d$，实测关节 $q_m$，名义 FK 为 $\hat p(q)$，真实位置为 $p_{\rm true}(q)$。严格加减同一项可得
\begin{align}
p_{\rm true}(q_m)-p_d={}&
\underbrace{p_{\rm true}(q_m)-\hat p(q_m)}_{\text{几何 安装和柔性误差}}\\
&+\underbrace{\hat p(q_m)-\hat p(q_d)}_{\text{关节跟踪引起的任务误差}}\\
&+\underbrace{\hat p(q_d)-p_d}_{\text{规划或逆解残差}}.
\end{align}
这三项应在相同坐标系和相同时刻比较。外部测量不准还会增加测量误差。把真实末端位置再次求一次 IK 后与 $q_d$ 比较，会受到冗余和多解影响，通常不是清楚的分解办法。

\begin{examplebox}[关节误差怎样变成毫米偏差]
在 $q_d=(0,\pi/2)$ 附近，测得 $\delta q=(0.01,-0.006)\,\mathrm{rad}$。
一阶近似给出
\[
\delta p\approx J_A\delta q=(-0.002,0.005)\,\mathrm m.
\]
用精确 FK 计算为约 $(-0.002025,0.004996)\,\mathrm m$，与一阶近似很接近。若外部相机测到偏差约 $(0.008,0.005)\,\mathrm m$，多出来的约 $10\,\mathrm{mm}$ 横向差不能再全部归因于这组关节跟踪误差，应检查几何、基座、TCP 或外部测量。
\end{examplebox}

\section{低阻抗回正偏差怎样从残差反推}

一维静态阻抗的核心关系为 $Ke=F_{\rm residual}$，其中残差的符号由统一动力学约定决定，可能包含重力估计误差、摩擦、线缆和外力。在 $K=100\,\mathrm{N/m}$ 时，$1\,\mathrm N$ 残差就对应 $1\,\mathrm{cm}$ 偏差。把允许偏差收紧到 $1\,\mathrm{mm}$，而刚度仍保持 $100\,\mathrm{N/m}$，则合成残差需要约束到 $0.1\,\mathrm N$ 量级。

这说明“软”和“准”可以通过更准的模型、更低摩擦或适当外层修正同时改善，但不能只把刚度调小然后要求相同扰动下误差不变。若增加积分，要重新考虑接触时积分积累和模式切换；若增加摩擦补偿，要考虑零速符号和过补偿能量注入。

从相反方向接近同一目标，可帮助区分方向相关和方向无关部分。例如两次静态位置误差为 $e_+=12\,\mathrm{mm}$、$e_-=4\,\mathrm{mm}$，同一方向刚度 $K=100\,\mathrm{N/m}$，对应残差 $1.2$ 和 $0.4\,\mathrm N$。平均值 $0.8\,\mathrm N$ 是方向无关部分的候选，半差 $0.4\,\mathrm N$ 是方向相关部分的候选。只有姿态、温度和其它条件近似相同，且摩擦近似对称时，这种分解才有解释力；平均值不能唯一证明是重力，半差也不能唯一证明是库仑摩擦。

\section{为什么不能凭方向说只是哪一个关节的问题}

原稿用“左右主要靠 J1，所以左右准”解释一个现象，这可以作为待验证假设，但不能作为任意构型下的结论。某方向末端运动需要哪些关节，由当前雅可比和约束共同决定；不同构型、姿态任务和零空间策略会改变分配。

更有用的做法是记录当前 $J$、关节力矩残差、任务误差、饱和状态及缩放后的奇异值。给定小位移方向 $d$，求受限的 $\delta q$，检查实际参与的关节。再看这些关节的重力残差、低速摩擦和力矩余量，而不是先按方向给关节定责。

\section{七轴的剩余自由度要服务于任务}

对满秩六维末端任务，七关节局部有一维运动学零空间；若只控制三维位置，局部可剩四维。剩余维度是 $n-\operatorname{rank}(J)$，不是永远固定为一。零空间可以帮助保持肘姿、远离限位、减少相机遮挡、降低重力力矩或避开线缆扭转。

“让肘部自由拖动”和“让肘部自动回到固定姿态”对应不同的零空间势能。较小或零的零空间刚度让肘部更自由，适当阻尼耗散运动；非零刚度会产生回正倾向。加上动态一致投影还要满足模型和秩条件，实际力矩限幅或阻尼伪逆会产生主任务耦合。第16章给出了这些条件下的完整公式和数值反例。

对相机云台，可以把任务分为过渡定位、离线运镜、在线跟随和拖动录制。它们共享模型与状态，但参考生成器和控制目标不同。画面平稳要同时检查工具姿态误差、角速度谱、相机安装、底座振动和时间同步，不能只看 TCP 平移的 RMS。

\section{一条能解释问题的学习与实施路线}

第一阶段只用两连杆完成 FK、IK、雅可比、动力学和功率核对。第二阶段做本章统一轨迹，区分纯反馈、模型前馈和执行器延迟。第三阶段通过独立训练与验证数据完成辨识，观察误差如何影响控制。第四阶段加入笛卡尔阻抗和外力观测，比较静态柔顺与动态稳定。第五阶段把同样的数据结构、单位和诊断量迁移到七轴模型，再分析冗余和真实工具。

每一阶段的交付证据应与问题对应：运动学需要位置和雅可比核对；辨识需要训练与独立验证残差；控制需要参考、测量、力矩和饱和曲线；实时实现需要数据年龄和耗时分布。不要把原稿里某个教学项目的“误差小于百分之五”“一毫秒循环”等数字直接当作所有真实机械臂的通用验收标准。具体指标应来自任务要求、传感器能力和硬件限制。

\section{计算练习与解答}

\subsection*{题一 底座误差}
相机到支撑转轴距离 $0.6\,\mathrm m$，允许支撑导致的位移误差 $0.5\,\mathrm{mm}$，小角度近似下允许角度是多少？

解：$\theta\le0.0005/0.6=0.00083333\,\mathrm{rad}=0.04775^\circ$。这是由支撑单独占用全部误差预算时的上界；其它误差还需分配预算。

\subsection*{题二 低刚度残差预算}
要求 $K=80\,\mathrm{N/m}$ 下静态误差不超过 $5\,\mathrm{mm}$，残余力应低于多少？

解：$|F_r|\le K|e|=80\times0.005=0.4\,\mathrm N$。这包括所有未被消除的等效残差，不仅是摩擦模型误差。

\subsection*{题三 七轴冗余维度}
七轴机械臂，某姿态下三维位置雅可比秩为三，六维几何雅可比秩为五。两种任务分别有多少维零空间？

解：位置任务为 $7-3=4$ 维；六维几何任务为 $7-5=2$ 维。但后者丢失了一个任务方向，不能把奇异点多出来的零空间当成获得了额外完整任务能力。

\section*{本章与原稿的对应}
本章整合 [A，第74--79页]、[B，第120--127页]，增加实际执行的统一 2R 数值实验、世界到相机误差分解及底座刚度预算。原稿中的 NERO 现场描述保留为诊断背景，所有新数值都明确为教学计算。
